\section{Discussion}

The discussion will focus on the difference between portability and measurement validity. Wearable EEG can reduce setup burden and enable more natural interaction, but these advantages increase the importance of transparent timing, movement context, artifact control, missing-data reporting, and device-specific limitations.

A major interpretive boundary is expected between studies showing group-level EEG differences and studies demonstrating reliable within-person event-locked measurement. The latter is the evidence required before using wearable EEG as a mechanistic companion measure in an adaptive digital intervention.

The review will also identify whether external reference validation, repeated-session reproducibility, and open preprocessing pipelines are common enough to support cross-study comparison.


\subsection{Implications of preliminary evidence mapping}

The supplied preliminary SciSpace synthesis suggests five testable reporting priorities: verified event-marker synchronization, quantified data yield and artifact handling, test--retest reliability, reference-device agreement, and recoverable preprocessing specifications. The apparent prevalence of deficiencies must not be interpreted before eligibility and full-text verification because the automated corpus includes secondary reviews, missing full texts and studies outside the primary ADHD use case. An embedded 2026 full-text pediatric ASD neurofeedback experiment provides contextual low-density acquisition detail but its principal pre/post resting-state network endpoint does not itself establish task-event synchronization or diagnostic utility in ADHD.

This distinction matters when designing experimental work: portable acquisition is a feasibility property; agreement of measured signals with a reference and reliable event timing are technical validity properties; individual-level inference and treatment effects require additional independent evidence. The review will evaluate each separately without treating a successful BCI feedback loop or between-group classifier as validation of a clinical biomarker.
